% Source-aligned English for printed p.213.
\clearpage
\SourcePage{213}

\LayerTarget{tohoku:proof:lemma-5.5.5}
\noindent $U_{s,\mathbf g}\cap f^{-1}(W^s_y)$ is zero. Such a $W^s_y$ exists
by the hypothesis on $P$, by the fact that there are only finitely many
systems $\mathbf g$ to consider, and because, as $W$ runs through a
fundamental system of neighborhoods of $y$, the sets
$U_{s,g}\cap f^{-1}(W)$ run through a fundamental system of
neighborhoods of the finite set of $x\in U_{s,\mathbf g}$ that project to $y$.
Since $Y$ is paracompact, a well-known lemma in Čech theory now gives a
covering ${\mathbf V}'=(V'_j)_{j\in J}$ of $Y$ and a map
$\varphi':J\to I$ such that $V'_j\subset U'_{\varphi'(j)}$ for every $j\in J$,
and such that for every $n$-simplex $t$ of $J$, $V'_t$ is contained in at
least one of the sets $W_y^{\varphi'(t)}$. Set
$V_j=U_{\varphi'(j)}\cap f^{-1}(V'_j)$ and
${\mathbf V}=(g\cdot V_j)_{(g,j)\in G\times J}$, and consider the map
$\varphi:G\times J\to G\times I$ defined by $\varphi'$. It is immediate that
$({\mathbf V},\varphi)$ satisfies the required conditions and, in particular,
$\varphi(F)=0$, whence $\overline\varphi(f^n)=0$.

The sheaf associated with the presheaf $H^q(A)$, for $q>0$, is zero (Lemma
3.8.2). Lemma 5.5.5 and the spectral sequence with initial term (5.5.15) then
give:

\LayerTarget{tohoku:theorem:5.5.6}
\textsc{Theorem 5.5.6.}\quad \textit{If $G$ is a discrete group of
homeomorphisms of $X$, and if $Y=X/G$ is paracompact, then the homomorphisms}

\LayerTarget{tohoku:equation:5.5.14}
\[
\tag{5.5.14}
\check H^n(X;G,A)\to H^n(X;G,A)
\]
\textit{are isomorphisms for every abelian $G$-sheaf $A$.}

\LayerTarget{tohoku:remark:2-low-degree-cech-consequences}
\textsc{Remarks II.}\quad From the spectral sequence with initial
term (5.5.15), without any hypothesis on $X,G,A$, one obtains

\LayerTarget{tohoku:equation:cech-h1-identification}
\[
\check H^1(X;G,A)=H^1(X,G,A)
\]
\LayerTarget{tohoku:comparison:cech-h2-injection}
 (and $\check H^2(X;G,A)\to H^2(X;G,A)$ is injective, since one
checks that $E_2^{0,n}=0$ for every $n>0$).
\LayerTarget{tohoku:interpretation:g-a-bundles}
Now $\check H^1(X;G,A)$ has a geometric interpretation, as one checks
easily, as the set of classes of $G$-$A$-bundles on $X$ (i.e. bundles on $X$
``with structural sheaf $A$'' [11], on which $G$ acts compatibly with this
structure). Moreover, $\check H^1(X;G,A)$ is also defined, and has the same
geometric interpretation, when $A$ is a $G$-sheaf of groups that need not be
commutative. The developments of [11], in particular Chapter V, can then be
carried over almost word for word to the more general context of $G$-bundles.
It would moreover be appropriate to give a treatment of ``noncommutative
Homological Algebra'' in the context of functors and categories, encompassing
both this theory of bundles and the algebraic mechanism of group-extension
questions (for Lie groups and others), as developed, for example, in papers
by Hochschild and A. Shapiro [12, 13, 19].

\LayerTarget{tohoku:section:5.6-equivariant-ext}
\textbf{5.6.\quad The $\operatorname{Ext}^n_{{\mathbf O},G}(X;A,B)$.}
In this subsection, ${\mathbf O}$ denotes a fixed $G$-sheaf of rings, and we consider
the abelian category ${\mathbf C}^{\mathrm O(G)}$ of $G$-${\mathbf O}$-modules (cf. 5.1).
\LayerTarget{tohoku:definition:equivariant-o-homomorphisms}
The group of $G$-${\mathbf O}$-homomorphisms from a $G$-${\mathbf O}$-module $A$ to another one
$B$ will be denoted $\operatorname{Hom}_{{\mathbf O},G}(A,B)$, or
$\operatorname{Hom}_{{\mathbf O},G}(X;A,B)$ if one wishes to make explicit

\TohokuContinuation{214}
